% ---------------------------------------------------------------------
%  APÊNDICE A — INJETOR DE VENTO
%  Movido dos Caps. 2 e 3 na revisão de 3/11 (item 3.7): nenhum ensaio de
%  atitude usa o vento, que no SITL quase não gera momento.
% ---------------------------------------------------------------------
\chapter{Injetor de vento}
\label{ap:vento}

O injetor de perturbações (Seção~\ref{sec:injetor}) também gera vento. No
modelo de multirrotor do SITL, o vento uniforme age como força no centro de
massa e quase não produz momento de atitude
(Seções~\ref{sec:sitl} e~\ref{sec:res-vento}); por isso, os ensaios de
rejeição de perturbação usam o torque, e o vento entra só no ensaio de rajada
da Seção~\ref{sec:res-vento}. Este apêndice reúne os modelos de vento, a sua
implementação no injetor e os ajustes do modelo de vento do SITL.

\section{Modelos de vento}
\label{ap:vento-modelos}

O vento que atua sobre a aeronave é usualmente decomposto em uma parcela
constante e uma parcela de rajadas, modelada como processo
estocástico~\cite{beard2012}; a rajada discreta da norma
MIL-F-8785C~\cite{mil8785c} é acrescentada aqui como terceira parcela. O
\emph{vento médio} é um vetor constante no referencial inercial; exige do
controlador ação integral ou compensação explícita. A \emph{rajada} discreta
é modelada por um perfil do tipo ``1~$-$~cosseno'', adaptado da
norma~\cite{mil8785c},
\begin{equation}
v_g(t) = \frac{A_g}{2}\left[1 - \cos\left(\frac{2\pi\,(t - t_0)}{T_g}\right)\right],
\qquad t_0 \le t \le t_0 + T_g,
\label{eq:rajada}
\end{equation}
de amplitude $A_g$ e duração $T_g$ conhecidas. Na norma, a rajada sobe ao
longo de uma distância $d_m$ e se mantém; aqui ela sobe em $T_g/2$ e volta a
zero, o que equivale a $d_m = V\,T_g/2$ para uma aeronave à velocidade $V$.
A \emph{turbulência} é um
processo estocástico obtido pela filtragem de ruído branco por funções de
conformação espectral. No modelo de Dryden, as densidades espectrais das
componentes longitudinal e lateral, em função da frequência espacial
$\Omega$, são~\cite[seção~3.7.1.2]{mil8785c}
\begin{equation}
\Phi_u(\Omega) = \sigma_u^2\,\frac{2L_u}{\pi}\,\frac{1}{1 + (L_u\Omega)^2},
\qquad
\Phi_v(\Omega) = \sigma_v^2\,\frac{L_v}{\pi}\,
\frac{1 + 3(L_v\Omega)^2}{\left[1 + (L_v\Omega)^2\right]^2},
\label{eq:dryden}
\end{equation}
em que $\sigma$ é a intensidade e $L$ a escala de comprimento da turbulência,
tabeladas em função da altitude e da severidade. Para uma aeronave à
velocidade $V$, a frequência temporal é $\omega = \Omega V$. A componente
longitudinal é um processo de Gauss-Markov de primeira ordem, de constante de
tempo $L_u/V$, discretizado exatamente pela
Equação~\eqref{eq:gauss-markov-discreto}, com $\tau = L_u/V$ e
$\sigma = \sigma_u$~\cite{gillespie1996}. A
componente lateral exige um filtro de segunda ordem~\cite{beard2012},
discretizado pela transformação de Tustin~\cite{tustin1947}.

\section{Implementação no injetor}
\label{ap:vento-injetor}

Em vento, o injetor gera o vento médio, a rajada do tipo
``1~$-$~cosseno''~\eqref{eq:rajada} e a turbulência de Dryden, com a
componente longitudinal pela discretização
exata~\eqref{eq:gauss-markov-discreto} e a lateral pelo filtro de segunda
ordem discretizado pela transformação de Tustin~\cite{tustin1947,franklin1998};
as duas normalizações foram conferidas numericamente, e ambas entregam o
desvio-padrão especificado. Como a aeronave paira, quem desloca o campo de
turbulência em relação a ela é o vento médio, e é a sua velocidade, com piso
de \SI{1}{\metre\per\second}, que o injetor usa no lugar de $V$; a escala
lateral é metade da longitudinal. Nenhum ensaio deste trabalho usa a
turbulência; o único ensaio de vento é a rajada da Seção~\ref{sec:res-vento}.

Três características do modelo de vento do SITL obrigaram a ajustar os seus
parâmetros, e o procedimento os fixa em todos os ensaios. Primeiro, o perfil
vertical de vento padrão (\codigo{SIM\_WIND\_T} = 0) reduz o vento perto do
solo segundo uma lei de raiz quadrada, e a \SI{5}{\metre} de altura chegariam
apenas 29\,\% do vento comandado; os ensaios usam o perfil sem redução.
Segundo, o vento passa por um filtro passa-baixas de primeira ordem cuja
constante de tempo (\codigo{SIM\_WIND\_TC}) vale \SI{5}{\second} por padrão,
o que apagaria qualquer rajada; os ensaios usam o piso imposto pelo firmware,
\SI{0.1}{\second}, e mesmo assim o vento fica atenuado acima de cerca de
\SI{1.6}{\hertz}. Terceiro, o SITL ignora o vento
nos primeiros \SI{5}{\second} de simulação; como o perfil só é disparado
depois da decolagem, mais de \SI{70}{\second} após a inicialização, essa
folga é sempre respeitada. A turbulência nativa do simulador fica desligada,
e a única turbulência é a do injetor.
